\chapter{Запись и чтение}
% полная версия: chapters/09_acet.tex

У памяти в мозге есть неудобная особенность: хранят воспоминания те же связи, по которым идёт работа. Нет отдельного жёсткого диска. Каждый раз, когда сеть вспоминает, она пользуется своими связями, --- и каждый раз, когда учится, она их меняет. Как не испортить старое, записывая новое? В компьютере на это есть отдельный сигнал --- «разрешить запись». В мозге его роль, похоже, играет радиостанция \nt{ACET} --- ацетилхолин.

\section*{Два режима}

Представьте сеть, которая хранит несколько картинок: покажи ей половину --- она дорисует остальное. Это чтение. Теперь покажите ей новую картинку, похожую на старую. Если сеть будет по привычке дорисовывать, она увидит не новое, а старое --- и выучит ошибку. Чтобы записать новое, сеть должна на время перестать доверять своей памяти и смотреть на мир.

\pencilfig{9}{\pencilart{9}}{Один провод переключает сеть: записывать новое или вспоминать старое.}

Ацетилхолин делает ровно это, и все три его эффекта измерены: он ослабляет внутренние связи сети, через которые она «вспоминает», усиливает сигналы, приходящие от органов чувств, и облегчает изменение связей. Много ацетилхолина --- режим записи: сеть слушает мир и учится. Мало --- режим чтения: сеть опирается на память и дорисовывает.

В нашей модели нет такого уровня, при котором хорошо шло бы и то, и другое. Для записи внутренние связи нужно приглушить, а для чтения они нужны в полную силу. Поэтому переключатель необходим.

\section*{Сколько доверять глазам}

Инженеры давно решили похожую задачу. Когда корабль идёт по курсу, у штурмана есть расчёт, где корабль должен быть, и есть новые, но неточные наблюдения. Сколько доверять каждому? Оптимальный фильтр отвечает: доверяй наблюдениям тем больше, чем меньше уверен в расчёте. И в этом же рецепте есть второе: чем меньше уверенность, тем быстрее надо переучиваться. Одна ручка управляет сразу и весом входа, и скоростью обучения. Это ровно то, что делает ацетилхолин. Отсюда гипотеза: он сообщает сети, насколько неуверенна её собственная модель мира.

\pencilfig{124}{%
% optimal blending: the less sure the model, the more weight on the observation (and the faster it relearns)
\foreach \dx/\wm/\we/\ttop/\bbot in {0/2.4pt/0.6pt/{модель уверена}/{ацетилхолина мало}, 4.3/0.6pt/2.4pt/{модель не уверена}/{ацетилхолина много}}{
  \begin{scope}[xshift=\dx cm]
  \draw[penlite=pcViolet, wash=pcViolet, rounded corners=3pt] (0,0.95) rectangle (1.3,1.45); \node[lbl, text=black] at (0.65,1.2) {память};
  \draw[penlite=pcBlue, wash=pcBlue, rounded corners=3pt] (0,-0.25) rectangle (1.3,0.25); \node[lbl, text=black] at (0.65,0) {глаза};
  \draw[dotpen=pcAmber, wash=pcAmber] (2.95,0.6) circle (0.55); \node[lbl, text=black] at (2.95,0.6) {оценка};
  \draw[pcViolet, line width=\wm, line cap=round, -{Stealth[length=5pt]}] (1.35,1.2) -- (2.4,0.8);
  \draw[pcBlue, line width=\we, line cap=round, -{Stealth[length=5pt]}] (1.35,0) -- (2.4,0.4);
  \node[font=\small\itshape, text=pcGray, anchor=south] at (1.55,1.6) {\ttop};
  \node[lbl, anchor=north] at (1.55,-0.4) {\bbot};
  \end{scope}}
}{Сколько доверять глазам. Если сеть уверена в своей модели мира, оценка почти вся из памяти. Если не уверена --- из наблюдений, и тогда же сеть быстрее переучивается. По гипотезе, эту ручку крутит ацетилхолин.}

\section*{Разделение труда}

Вспомните прошлую главу. Когда мир внезапно меняется, норадреналин, по гипотезе, замечает неожиданность. А ацетилхолин держит сеть в режиме записи, пока новая обстановка не выучена. Один сообщает «что-то изменилось», другой --- «пиши».

\section*{Ночная перезапись}

Во время глубокого сна уровень ацетилхолина низок. Сеть отключена от мира и находится в режиме чтения. Именно в это время в ней «проигрываются» дневные события --- это измерено. Что при этом выученное за день переписывается в долговременную память, --- правдоподобная гипотеза, к которой мы вернёмся в главе о сне.

\fullversion{9}
